\scanpage{9}
\deleted{À tout ``cône'' \unsure{stable|stabl}}

Si nous nous intéressons aux structures d'effectivité \inserted{qui peuvent se} définir \deleted{par} en termes de \deleted{\unsure{certain}} une famille de $\underline{i}$ \deleted{\unsure{pour}}comme intersection des $P(i)$, \deleted{\unsure{on a}} la donnée d'une telle structure équivaut à celle d'un ``cône'' de $\Gamma'$, invariant par $\pi_1$, contenu \inserted{``\uline{strictement}''} dans celui défini par $\varepsilon$,
% slanted
\margin{\unsure{``''|!!!} i.e. $\varepsilon$ est \unsure{``}\uline{intérieur} au cône polaire !}
contenant la classe de $i$, satisfaisant à une condition de \deleted{\unsure{non} aplatissement} saturation que nous ne précisons pas ici (est-elle automatiquement satisfaite ?).

Le cas le plus intéressant est celui \deleted{\unsure{de}} \inserted{où} $M^+$ \deleted{\unsure{$Q$|$\otimes$}} est défini comme $P(i_1)$, où $i_1$ est comme au $3^{\circ}$. \deleted{\unsure{En particulier},} \inserted{[N.B.} le cône \uline{invariant} engendré contient bien $j$, et est contenu dans le $\frac12$ espace défini par $\varepsilon$, puisque $i_1\bar i_1=j$, $\varepsilon i_1=\mathrm{id}$]. On trouve ainsi une application
\[
I(G)\longrightarrow\deleted{\unsure{\mathrm{Pol}}}\mathrm{Eff}(G)
\]
% archive sheet number: 8
